\section{Introduction}
% Cold War as introduction to theories of IR
After the ending of the Cold War, a vaccuum in International Relations (IR) theory was needed to be filled, as the existing paradigms couldn't explain its sudden and peaceful ending. After almost half a century of geopolitical tensions, characterized by proxy wars, a nuclear arms race and technological rivalry, George H.W. Bush and Mikhail Gorbachev declared the end of the Cold War in Malta on December 3rd, 1989. Consequently, the power balance shifted, introducing a unipolarity of the United States and, importantly, increasing the influence of NGOs in the national affairs of states.
 
% Introduction Neorealism & Neoliberalism
% Evolution of Realism into Neorealism in the Latter half of the Cold War
Since the latter half of the Cold War, \textbf{Neorealism} is seen as one of the main pillars in International Relations theory. Kenneth Waltz formalized neorealism in his 1979 book \textit{Theory of International Relations}, building upon the basis of classical realism. Classical realism explains the tides of International Relations through the lens of human nature, making it subject to the ego and emotions of world leaders. Neorealism shifts the focus from human nature to international systemic structure, but agrees with classical realism that states act out of self-interest, be it for survival or power accumulation, and that relative gains in material capabilties are more important than absolute gains \cite{waltz1979, morgenthau1978}.

% Introduction Constructivism
In 1992, Wendt \cite{wendt1992} propagates a new claim, namely that of CONSTRUCTIVISM. He disputes the theory of Waltz by stating that states are not inherently focused on self-optimization, but that their identities and interests are formed through their interactions with others. In essence, states are not self-optimizing per se. Percieved semantic meaning of interactions determine identity and interest. States can become self-optimizing, but not because they are. Process determines the outcomes.

% Comparison Neorealism vs Constructivism

% State of debate today 


% Relevance to thesis


\subsection{Computational Mapping}

The thesis will try to map these theories to a computational domain.

\textbf{REINFORCEMENT LEARNING AGENTS} are pure utility optimisers. Their policies and actions are based upon pure game theoretic calculation that ensures their optimal long-term utility accumulation. If we set their utility function to optimising "power" within the system, they are essentially REALIST agents/states.

% TODO: [Misschien aanvullen met Reward is Enough inzichten?]
% Consider adding \cite{silver2021reward} reference

\textbf{LLMs} are semantic and perhaps complex stochastic parrots. They are trained upon the human corpus and semantic interpreters. They will be able to assign meaning to interactions and translate that computationally.

% TODO: [Inzichten GENERATIVE AGENTS + ]
% TODO: [Merk dat ik moeite heb met het mappen van constructivisme aan LLM agents]
% Questions to address:
% - How do LLMs construct narratives from interaction history?
% - Why is semantic interpretation different from pattern matching?
% - What does it mean for identity/interests to be "emergent" in LLMs vs fixed in RL?
